%=============================================================================
% SECTION 8B: GAUGE COUPLING VARIATION AND HIGH-ENERGY IMPLICATIONS
% Part V: Parameter-Free Predictions (continued)
% Extends the α-relations to all gauge couplings
%=============================================================================

\section{Gauge Coupling Variation and High-Energy Implications}\label{sec:gauge-coupling}

Section~\ref{sec:alpha} established that electromagnetic properties couple to the scalar field $\psi$ through $\kalpha = \alphafine^2/(2\pi)$. This section extends the framework to all Standard Model gauge couplings, derives the modified renormalization group equations, and explores consequences ranging from nuclear clock tests to grand unification.

%-----------------------------------------------------------------------------
\subsection{Universal Gauge-$\psi$ Coupling}\label{sec:gauge-universal}
%-----------------------------------------------------------------------------

\paragraph{Extension to all gauge sectors.}
The clock coupling $\kalpha = \alphafine^2/(2\pi)$ arises from the interaction between electromagnetic fields and the DFD optical metric. A parallel derivation for non-Abelian gauge fields yields the universal form:
\begin{equation}\label{eq:ki-universal}
\frac{\delta\alpha_i}{\alpha_i} = k_i\,\psifield, \qquad k_i = \frac{\alpha_i^2}{2\pi},
\end{equation}
where $\alpha_i = g_i^2/(4\pi)$ is the fine-structure constant for gauge group $i$.

\paragraph{Physical origin.}
The $\alpha_i^2$ dependence is characteristic of one-loop quantum corrections. The optical metric $\tilde{g}_{\mu\nu} = e^{2\psi}\eta_{\mu\nu}$ modifies gauge field propagators, and quantum corrections generate this dependence through loop diagrams. The gauge emergence framework (Appendix~\ref{app:gauge-emergence}) provides a deeper origin for these couplings through frame stiffness in the internal mode space.

\paragraph{The gauge hierarchy.}
At laboratory energies:
\begin{align}
\text{U(1)}_{\rm EM}:& \quad \alphafine \approx 1/137, \quad \kalpha \approx 8.5 \times 10^{-6}, \\
\text{SU(2)}_L:& \quad \alpha_w \approx 1/30, \quad k_w \approx 1.8 \times 10^{-4}, \\
\text{SU(3)}_c:& \quad \alpha_s \approx 0.118, \quad k_s \approx 2.2 \times 10^{-3}.
\end{align}

The strong force is most sensitive to gravitational potential:
\begin{equation}\label{eq:ks-kalpha-ratio}
\frac{k_s}{\kalpha} = \frac{\alpha_s^2}{\alphafine^2} \approx 260.
\end{equation}

\begin{tcolorbox}[colback=green!5!white, colframe=green!50!black, title=The Gauge Coupling Hierarchy]
\textbf{Key result:} All gauge couplings shift with gravitational potential according to $\delta\alpha_i/\alpha_i = k_i\psi$ with $k_i = \alpha_i^2/(2\pi)$.

\textbf{Hierarchy:} $k_s : k_w : \kalpha \approx 260 : 20 : 1$

The strong force is $\sim 260\times$ more sensitive to $\psi$ than electromagnetism.
\end{tcolorbox}

%-----------------------------------------------------------------------------
\subsection{Connection to the $\beta$-Function}\label{sec:gauge-beta}
%-----------------------------------------------------------------------------

\paragraph{The one-loop $\beta$-function.}
The running of gauge couplings with energy scale $\mu$ is governed by:
\begin{equation}\label{eq:beta-standard}
\frac{d\alpha_i}{d\ln\mu} = \frac{b_i\alpha_i^2}{2\pi},
\end{equation}
where $b_i$ is the one-loop coefficient:
\begin{align}
b_1 &= +\frac{41}{10} \quad \text{(U(1)}_Y\text{)}, \\
b_2 &= -\frac{19}{6} \quad \text{(SU(2)}_L\text{)}, \\
b_3 &= -7 \quad \text{(SU(3)}_c\text{)}.
\end{align}

\paragraph{The remarkable connection.}
Comparing Eqs.~\eqref{eq:ki-universal} and~\eqref{eq:beta-standard}:
\begin{equation}\label{eq:ki-beta}
\boxed{k_i = \frac{\beta_i}{b_i}}
\end{equation}

\textbf{The $\psi$-gauge coupling equals the $\beta$-function divided by the group-theory coefficient.}

\paragraph{Physical interpretation.}
This reveals that gravitational potential acts as an effective shift in the renormalization scale. Gravity and RG flow are connected at all energy scales through $k_i = \alpha_i^2/(2\pi)$.

%-----------------------------------------------------------------------------
\subsection{Modified Renormalization Group Equations}\label{sec:gauge-rg}
%-----------------------------------------------------------------------------

In the presence of non-zero $\psi$, gauge couplings depend on both energy scale and gravitational potential:
\begin{equation}
\alpha_i(\mu, \psi) = \alpha_i(\mu, 0)\left(1 + \frac{\alpha_i^2}{2\pi}\psi\right).
\end{equation}

Taking the scale derivative at fixed $\psi$:
\begin{align}
\frac{d\alpha_i(\mu,\psi)}{d\ln\mu} &= \frac{d\alpha_i(\mu,0)}{d\ln\mu}
\left(1 + \frac{\alpha_i^2\psi}{2\pi}\right) 
+ \alpha_i \cdot \frac{2\alpha_i}{2\pi}\frac{d\alpha_i}{d\ln\mu}\psi.
\end{align}

\textbf{The modified $\beta$-function:}
\begin{equation}\label{eq:beta-modified}
\boxed{\frac{d\alpha_i}{d\ln\mu} = \frac{b_i\alpha_i^2}{2\pi}
\left[1 + \frac{3\alpha_i^2}{2\pi}\psi\right]}
\end{equation}

The $\psi$-correction is proportional to $\alpha_i^4$---a \textbf{two-loop-like gravitational correction} to the running.

\paragraph{Laboratory effects.}
For QCD near confinement ($\alpha_s \sim 1$):
\begin{equation}
\frac{\delta\beta_s}{\beta_s} \sim \frac{\alpha_s^2\psi}{2\pi} \sim 0.05\psi.
\end{equation}
In laboratory environments ($\psi \sim 10^{-9}$), this is $\sim 10^{-10}$---unmeasurable directly, but the $k_s$ coupling itself has dramatic consequences for nuclear physics.

%-----------------------------------------------------------------------------
\subsection{Asymptotic Freedom and UV Behavior}\label{sec:gauge-asymptotic}
%-----------------------------------------------------------------------------

\paragraph{QCD decoupling.}
QCD is asymptotically free: $\alpha_s(\mu) \to 0$ as $\mu \to \infty$. This implies:
\begin{equation}
k_s(\mu) = \frac{\alpha_s^2(\mu)}{2\pi} \to 0 \quad \text{as } \mu \to \infty.
\end{equation}

\textbf{The strong sector decouples from $\psi$ in the ultraviolet.}

\paragraph{Maximum sensitivity at confinement.}
Conversely, $k_s$ is \textit{maximal} at the confinement scale where $\alpha_s \sim 1$:
\begin{equation}
k_s^{\rm max} \sim \frac{1}{2\pi} \approx 0.16.
\end{equation}

This explains why nuclear physics provides the strongest low-energy probe of $\psi$-gauge coupling: the effective coupling $k_s$ peaks precisely at the energy scale relevant for nuclear binding.

\paragraph{QED behavior.}
QED is not asymptotically free; $\alphafine$ increases slowly with energy. The Landau pole occurs at $\mu \sim 10^{286}$ GeV, far above the Planck scale. For practical purposes, $\kalpha$ remains approximately constant.

%-----------------------------------------------------------------------------
\subsection{Nuclear Clock Prediction: Thorium-229}\label{sec:gauge-nuclear}
%-----------------------------------------------------------------------------

The $k_s/\kalpha \approx 260$ hierarchy, combined with the exponential sensitivity of QCD through dimensional transmutation, leads to dramatic predictions for nuclear transitions.

\paragraph{The thorium-229 isomer.}
$^{229}$Th has a nuclear isomer with uniquely low transition energy:
\begin{equation}
E_m = 8.338 \pm 0.024 \text{ eV}.
\end{equation}
This arises from near-cancellation between Coulomb ($\sim +300$ keV) and nuclear strong-force ($\sim -300$ keV) contributions, with a residual of only $\sim 8$ eV.

\paragraph{Sensitivity coefficients.}
The isomer energy depends on fundamental constants through:
\begin{equation}\label{eq:Em-sensitivity}
\frac{\delta E_m}{E_m} = K_\alpha\frac{\delta\alphafine}{\alphafine} + K_q\frac{\delta X_q}{X_q},
\end{equation}
where $X_q \equiv m_q/\Lambda_{\rm QCD}$ and from nuclear structure calculations:
\begin{equation}
K_\alpha \approx 10^4, \qquad K_q \approx -10^4.
\end{equation}

\paragraph{The $\Lambda_{\rm QCD}$ amplification.}
The QCD scale is determined by dimensional transmutation:
\begin{equation}
\Lambda_{\rm QCD} = \mu\exp\left(-\frac{2\pi}{|b_3|\alpha_s(\mu)}\right).
\end{equation}
Differentiating:
\begin{equation}\label{eq:LQCD-sensitivity}
\frac{\delta\Lambda_{\rm QCD}}{\Lambda_{\rm QCD}} = \frac{2\pi}{|b_3|\alpha_s^2}\,\delta\alpha_s = \frac{2\pi}{|b_3|\alpha_s}\,\frac{\delta\alpha_s}{\alpha_s} \approx 7.6\,\frac{\delta\alpha_s}{\alpha_s}.
\end{equation}

The factor $2\pi/(|b_3|\alpha_s) \approx 7.6$ represents the exponential amplification of relative coupling changes through dimensional transmutation.  (Note: the coefficient of the \emph{absolute} change $\delta\alpha_s$ is the larger number $2\pi/(|b_3|\alpha_s^2) \approx 64$; these two bookkeeping conventions must not be mixed.)

\paragraph{The DFD enhancement factor.}
Combining the above with $\delta X_q/X_q \approx -\delta\Lambda_{\rm QCD}/\Lambda_{\rm QCD}$ and using $\delta\alpha_s/\alpha_s = k_s\psi$:
\begin{align}
\frac{\delta E_m}{E_m} &= K_\alpha\kalpha\psi + K_q \times \frac{2\pi}{|b_3|\alpha_s}\,k_s\,\psi \nonumber\\
&= \left(10^4 \times 8.5 \times 10^{-6} - 7.6 \times 10^4 \times 2.2 \times 10^{-3}\right)\psi \nonumber\\
&\approx (0.085 - 167)\psi \approx -170\,\psi.
\end{align}

For comparison, an optical atomic clock has $\delta\nu_{\rm opt}/\nu_{\rm opt} \approx \psi$.

\begin{tcolorbox}[colback=red!5!white, colframe=red!60!black, title=Nuclear Clock Enhancement: Unscreened Gauge-Sector Estimate]
\begin{equation}\label{eq:enhancement}
\boxed{\mathcal{R} \equiv \frac{(\delta\nu/\nu)_{\rm Th\text{-}229}}{(\delta\nu/\nu)_{\rm optical}} \approx -170^{+300}_{-120}}
\end{equation}

\textbf{Caveat:} This is the \emph{unscreened} gauge-sector estimate.  The screened treatment in Sec.~\ref{sec:clocks-nuclear}, incorporating the $\mu_{\rm LPI}$ screening function and 2026 Th-229 reproducibility data, substantially reduces the expected amplitude and compresses the surviving annual signal window to 26\,Hz--$\mathcal{O}(1\,\mathrm{kHz})$.  The unscreened value above serves as the theoretical ceiling, not the experimental target.

\textbf{Physical origin:}
\begin{enumerate}
\item $k_s \gg \kalpha$: Strong force couples to $\psi$ more strongly
\item Dimensional transmutation: $\Lambda_{\rm QCD}$ exponentially sensitive to $\alpha_s$
\item Near-cancellation: 8 eV isomer is tiny residual of ${\sim}$MeV forces
\end{enumerate}
\end{tcolorbox}

\paragraph{Experimental test protocol.}
The following estimates use the \emph{unscreened} enhancement $|\mathcal{R}| \approx 170$.  The screened predictions, which are the operationally relevant ones for terrestrial experiments, are given in Sec.~\ref{sec:clocks-nuclear}.

\textbf{Height experiment (1 m separation), unscreened:}
\begin{align}
\text{GR:}& \quad \frac{\Delta(\nu_{\rm Th}/\nu_{\rm Sr})}{\nu_{\rm Th}/\nu_{\rm Sr}} = 0, \\
\text{DFD (unscreened):}& \quad \frac{\Delta(\nu_{\rm Th}/\nu_{\rm Sr})}{\nu_{\rm Th}/\nu_{\rm Sr}} \approx 1.8 \times 10^{-14}.
\end{align}

\textbf{Annual modulation (solar potential), unscreened:}
\begin{align}
\text{GR:}& \quad \left|\frac{\Delta(\nu_{\rm Th}/\nu_{\rm Sr})}{\nu_{\rm Th}/\nu_{\rm Sr}}\right|_{\rm annual} = 0, \\
\text{DFD (unscreened):}& \quad \left|\frac{\Delta(\nu_{\rm Th}/\nu_{\rm Sr})}{\nu_{\rm Th}/\nu_{\rm Sr}}\right|_{\rm annual} \approx 5 \times 10^{-8}.
\end{align}

\paragraph{Timeline.}
$^{229}$Th nuclear clocks are under active development:
\begin{itemize}
\item 2024: First laser excitation of nuclear transition demonstrated
\item 2026--27: First-generation nuclear clocks at $\sim 10^{-12}$ precision
\item 2028--30: Improved precision to $\sim 10^{-15}$
\end{itemize}

\textbf{The DFD prediction is testable within 2--3 years.}

%-----------------------------------------------------------------------------
\subsection{Cosmological $\alpha(z)$ Variation}\label{sec:gauge-cosmology}
%-----------------------------------------------------------------------------

If the cosmological gravitational potential $\psi$ evolves with redshift, then $\alpha$ evolves accordingly.

\paragraph{Cosmological potential.}
In DFD, the cosmological scalar field tracks the matter density:
\begin{equation}\label{eq:psi-cosmo}
\psi(z) = \frac{\xi_{\rm LPI}^{\rm res}}{2}\Omega_m(z),
\end{equation}
where $\xi_{\rm LPI}^{\rm res}$ is the residual screened cavity/clock coupling scale discussed in Sec.~\ref{sec:cavity} and
\begin{equation}
\Omega_m(z) = \frac{\Omega_{m,0}(1+z)^3}{\Omega_{m,0}(1+z)^3 + \Omega_\Lambda}.
\end{equation}

\paragraph{The $\alpha(z)$ prediction.}
Combining with $\kalpha = \alphafine^2/(2\pi)$:
\begin{equation}\label{eq:alpha-z}
\frac{\Delta\alphafine}{\alphafine}(z) = \kalpha[\psi(z) - \psi_0] = \frac{\xi_{\rm LPI}^{\rm res}\alphafine^2}{4\pi}\left[\Omega_m(z) - \Omega_{m,0}\right].
\end{equation}

For illustrative plotting one may temporarily set $\xi_{\rm LPI}^{\rm res}=1$, but the corrected cavity sector indicates that the physically relevant value is a much smaller screened residual:
\begin{equation}
\frac{\Delta\alphafine}{\alphafine}(z) \approx 7 \times 10^{-6} \times \left[\Omega_m(z) - 0.31\right].
\end{equation}

\paragraph{Numerical predictions.}
\begin{center}
\renewcommand{\arraystretch}{1.2}
\begin{tabular}{lccc}
\toprule
Epoch & Redshift & $\Omega_m(z)$ & $\Delta\alpha/\alpha$ (DFD) \\
\midrule
Quasars & 2 & 0.91 & $+4 \times 10^{-6}$ \\
CMB & 1100 & 1.00 & $+5 \times 10^{-6}$ \\
BBN & $10^9$ & 1.00 & $+5 \times 10^{-6}$ \\
\bottomrule
\end{tabular}
\end{center}

\paragraph{Comparison with observational bounds.}

\textbf{Laboratory input.}
In DFD the cosmological $\alpha$-variation is controlled by the same residual LPI scale $\xi_{\rm LPI}^{\rm res}$ discussed for cavity--atom tests (Sec.~\ref{sec:cavity}). We treat $\xi_{\rm LPI}^{\rm res}$ as an experimentally determined input, not a cosmology fit parameter. Cosmological bounds therefore constrain the laboratory value of this residual scale.

\begin{table}[h]
\centering
\caption{Observational probes of fine-structure constant variation.}\label{tab:alpha-variation}
\begin{ruledtabular}
\scriptsize
\begin{tabular}{lccl}
\textbf{Probe} & $z$ & \textbf{DFD pred.} & \textbf{Observed} \\
\hline
ESPRESSO & $0.6$--$2.4$ & $+4\xi_{\rm LPI}^{\rm res}$ ppm & $(-0.5 \pm 0.6)$ ppm \\
Quasar dipole & $1$--$3$ & --- & $\sim10$ ppm \\
CMB & 1100 & $+5\xi_{\rm LPI}^{\rm res}$ ppm & $<2000$ ppm \\
BBN & $10^9$ & $+5\xi_{\rm LPI}^{\rm res}$ ppm & $<20000$ ppm \\
\end{tabular}
\end{ruledtabular}
\vspace{1mm}
{\scriptsize References: ESPRESSO~\cite{Murphy2022ESPRESSO}; dipole~\cite{Webb2011,Whitmore2015}; CMB~\cite{Planck2018VI}; BBN~\cite{Cyburt2016}.}
\end{table}

Using the conservative ppm-level quasar constraints, the scaling $\Delta\alpha/\alpha \sim (4\times 10^{-6})\,\xi_{\rm LPI}^{\rm res}$ implies that a genuinely order-unity cosmological residual would already be uncomfortable. The corrected cavity sector therefore pushes this subsection into the category of a conditional screen/coupling dictionary rather than a settled laboratory-normalized result.

\textbf{Status:}
\begin{itemize}
\item BBN and CMB: Satisfied for $\xi_{\rm LPI}^{\rm res} \leq 1$ with $>100\times$ margin.
\item Quasars: For $\xi_{\rm LPI}^{\rm res}$ of order unity, bounds become constraining. Current quasar systematics are debated~\cite{Whitmore2015}.
\item The cosmological prediction is only as clean as the laboratory determination of the residual coupling scale; with the cavity correction, this subsection should be read as conditional rather than closed.
\end{itemize}

\paragraph{Distinctive signatures.}
DFD predicts specific features distinguishing it from other varying-$\alpha$ models:
\begin{enumerate}
\item \textbf{Functional form:} $\Delta\alpha/\alpha$ tracks $\Omega_m(z)$, flat at high $z$ and falling steeply for $z < 1$
\item \textbf{Sign:} $\Delta\alpha/\alpha > 0$ (larger $\alpha$ in the past)
\item \textbf{Spatial correlation:} $\Delta\alpha/\alpha$ should correlate with local matter density
\end{enumerate}

\paragraph{Future tests.}
The ELT/ANDES spectrograph will achieve $\sigma(\Delta\alpha/\alpha) \sim 10^{-7}$ per quasar system, tightening constraints on the residual cosmological coupling scale $\xi_{\rm LPI}^{\rm res}$ and potentially detecting a ppm-level signal if that screened residual lies near the upper end allowed by the clock sector.

%-----------------------------------------------------------------------------
\subsection{Grand Unification}\label{sec:gauge-gut}
%-----------------------------------------------------------------------------

\paragraph{Standard unification picture.}
The SM gauge couplings approximately unify at $M_{\rm GUT} \sim 10^{15-16}$ GeV, but with a mismatch of $\sim 3$--5\%.

\paragraph{DFD corrections.}
Couplings measured today include $\psi$-corrections from cosmological evolution:
\begin{equation}
\alpha_i^{\rm today} = \alpha_i^{\rm GUT}\left(1 + k_i^{\rm low}\Delta\psi\right),
\end{equation}
where $\Delta\psi = \psi_{\rm today} - \psi_{\rm GUT}$ and $|\Delta\psi| \sim 1$.

\paragraph{Differential corrections.}
\begin{align}
\frac{\delta\alpha_1}{\alpha_1} &\approx 5 \times 10^{-5}, \\
\frac{\delta\alpha_2}{\alpha_2} &\approx 2 \times 10^{-4}, \\
\frac{\delta\alpha_3}{\alpha_3} &\approx 2 \times 10^{-3}.
\end{align}

\paragraph{Effect on unification.}
The relative shift in the unification condition:
\begin{equation}
\frac{\delta(\alpha_3 - \alpha_1)}{\alpha_{\rm GUT}} \sim (k_3 - k_1)\Delta\psi \sim 2 \times 10^{-3}.
\end{equation}

\textbf{DFD predicts a $\sim 0.2\%$ shift in gauge coupling unification.}

Since $k_3 > k_2 > k_1$ and $\Delta\psi > 0$ (larger $\psi$ in the past), the correction slightly \textit{worsens} unification---about 5\% of the total SM mismatch. This is smaller than current theoretical uncertainties but represents a definite prediction.

%-----------------------------------------------------------------------------
\subsection{Vacuum Energy Feedback}\label{sec:gauge-cc}
%-----------------------------------------------------------------------------

The $\psi$-gauge coupling creates a feedback loop connecting vacuum energy, gravitational potential, and gauge couplings:
\begin{center}
$\rho_{\rm vac} \xrightarrow{\text{source}} \psi \xrightarrow{\text{shift}} \alpha_i \xrightarrow{\text{loops}} \rho_{\rm vac}$
\end{center}

\paragraph{Self-consistency condition.}
Let $\psi = F(\rho_{\rm vac})$ be the sourcing relation and $\rho_{\rm vac} = G(\alpha_i(\psi))$ be the loop contribution. Fixed points satisfy $\psi^* = \Phi(\psi^*)$.

\paragraph{Stability analysis.}
Linearizing around $\psi = 0$:
\begin{equation}
\psi^* = \frac{\psi_0}{1 - \lambda},
\end{equation}
where:
\begin{equation}
\lambda \sim \frac{M_P^4}{\rho_c} \times \frac{\alpha^3}{128\pi^3} \sim 10^{113}.
\end{equation}

\textbf{The feedback is violently unstable: $\lambda \sim 10^{113} \gg 1$.}

\paragraph{Interpretation.}
The enormous value of $\lambda$ means small perturbations in $\psi$ grow by a factor of $\sim 10^{113}$ per iteration. Possible interpretations:
\begin{enumerate}
\item Self-tuning to $\psi = 0$ as the only stable fixed point
\item UV cutoff constraint: proper UV completion must regulate this feedback
\item New physics required for stabilization
\end{enumerate}

\textbf{Constraint on UV completion:} Any UV completion of DFD must make the $\psi$-vacuum energy feedback loop stable. Note that the cosmological constant problem is solved separately by topology: $(H_0/M_P)^2 = \alpha^{57}$ (Section~\ref{sec:G-derivation}). This feedback loop concern is about UV stability, not the $\Lambda$ value.

%-----------------------------------------------------------------------------
\subsection{Summary of Falsifiable Predictions}\label{sec:gauge-predictions}
%-----------------------------------------------------------------------------

\begin{table}[h]
\centering
\caption{Tier 1: Nuclear clock tests (unscreened gauge-sector estimates; see Sec.~\ref{sec:clocks-nuclear} for screened predictions)}
\label{tab:tier1-predictions}
\renewcommand{\arraystretch}{1.2}
\begin{tabular}{lccc}
\toprule
Observable & GR & DFD (unscreened) & Timeline \\
\midrule
Th/Sr ratio (1m height) & 0 & $1.8 \times 10^{-14}$ & 2026--27 \\
Th/Sr annual modulation & 0 & $5 \times 10^{-8}$ & 2026--27 \\
Nuclear vs optical sign & Same & Opposite & 2026--27 \\
\bottomrule
\end{tabular}
\end{table}

\textbf{Note:} These are unscreened estimates.  The screened treatment in Sec.~\ref{sec:clocks-nuclear}, incorporating $\mu_{\rm LPI}$ screening and 2026 Ooi reproducibility data, compresses the surviving signal window to 26\,Hz--$\mathcal{O}(1\,\mathrm{kHz})$.  If the measured Th/Sr enhancement is consistent with unity at $5\sigma$ \emph{and} the cross-species atomic channels also show persistent nulls, the DFD gauge-sector coupling structure would be falsified.

\begin{table}[h]
\centering
\caption{Tier 2: Constraining medium-term tests}\label{tab:tier2-predictions}
\begin{ruledtabular}
\scriptsize
\begin{tabular}{lccc}
\textbf{Observable} & \textbf{DFD pred.} & \textbf{Current} & \textbf{Test} \\
\hline
$\Delta\alpha/\alpha$ ($z\sim 2$) & $\approx 4\xi_{\rm LPI}^{\rm res}$ ppm & ppm-level & ELT \\
$\alpha(z)$ shape & $\propto \Omega_m(z)$ & --- & ELT \\
Spatial $\alpha$ corr. & $\propto \delta_m$ & --- & ELT \\
\end{tabular}
\end{ruledtabular}
\end{table}

\begin{table}[h]
\centering
\caption{Tier 3: Theoretical consistency tests}\label{tab:tier3-predictions}
\renewcommand{\arraystretch}{1.2}
\small
\begin{tabular}{@{}lcc@{}}
\toprule
Quantity & DFD prediction & Status \\
\midrule
GUT shift & $\sim 0.2\%$ & Below precision \\
Modified $\beta$ & $\delta\beta \propto \alpha^4\psi$ & Unmeasurable \\
CC feedback & $\lambda \sim 10^{113}$ & Constrains UV \\
\bottomrule
\end{tabular}
\end{table}

\paragraph{Hierarchy of tests.}
\begin{enumerate}
\item \textbf{Nuclear clocks} test the core relation $k_i = \alpha_i^2/(2\pi)$. Confirmation validates the entire gauge-$\psi$ framework.
\item \textbf{Cosmological $\alpha(z)$} tests the $\psi$-cosmology connection, independent of nuclear physics uncertainties.
\item \textbf{GUT and CC constraints} test high-energy implications, relevant once Tiers 1--2 are confirmed.
\end{enumerate}

\begin{tcolorbox}[colback=blue!5!white, colframe=blue!50!black, title=Summary: Gauge Coupling Variation]
\textbf{Universal coupling:} $\delta\alpha_i/\alpha_i = k_i\psi$ with $k_i = \alpha_i^2/(2\pi)$

\textbf{Key insight:} $k_i = \beta_i/b_i$ --- gravity acts as effective RG scale shift

\textbf{Hierarchy:} $k_s : k_w : \kalpha \approx 260 : 20 : 1$

\textbf{Nuclear clock (unscreened):} $\mathcal{R} \approx -170$; screened predictions in Sec.~\ref{sec:clocks-nuclear}

\textbf{Cosmological $\alpha$:} $\Delta\alpha/\alpha \sim 5 \times 10^{-6}$ from BBN to today

\textbf{Falsification criteria:}
\begin{itemize}
\item Persistent nulls across all clock channels (same-ion, cross-species, nuclear) falsifies the gauge-sector framework
\item $\mathcal{R} \approx 1$ with high precision rules out DFD gauge coupling
\item $|\mathcal{R}| \sim 10^2$ with correct sign: strong confirmation
\end{itemize}
\end{tcolorbox}
